\usetikzlibrary{arrows.meta,
                positioning,
                shadows}

\usetikzlibrary{shapes.multipart,positioning}

\usetikzlibrary{shapes, arrows, arrows.meta, fit,backgrounds, positioning, calc}

\tikzstyle{process} = [rectangle, minimum width=2cm, minimum height=1cm, text centered, text width=2cm,  draw=black, fill=orange!30]
\tikzstyle{arrow} = [thick,->,>=stealth]
\tikzstyle{container1} = [draw, rectangle, inner sep=0.3cm, fill=gray!20]
\tikzset{
  mybackground/.style={execute at end picture={
      \begin{scope}[on background layer]
        \node[] at (current bounding box.north){\bottom{1cm} #1};
        \end{scope}
    }},
}

\tikzset{
    font=\sffamily,
    BLOCK/.style={
        draw,
        align=center,
        %
        draw=red!50,
        fill=red!20,
        rectangle split, 
        rectangle split horizontal,
        rectangle split parts=#1, 
    }
}
\tikzset{cross/.style={cross out, draw, 
         minimum size=2*(#1-\pgflinewidth), 
         inner sep=0pt, outer sep=0pt}}
\begin{figure}[t]
\centering  
		\resizebox{.8\columnwidth}{!}{  
        \begin{tikzpicture}[%
            every path/.style={thick,%
            },
                           dcs/.style = {double copy shadow, shadow xshift=4pt, shadow yshift=-4pt}
          ]
        \node [rectangle, minimum width=.5cm, minimum height=1cm, text centered, text width=1cm,  draw=black, fill=blue!30] (process3) at (-2, -1.5) {$\mathbf{S}_{[i-1]}$};

        \node [rectangle, minimum width=.5cm, minimum height=.75cm, text centered, text width=.75cm,  draw=black, fill=red!20] (hh1) at (-2, -6.5) {$\mathbf{S}_{[i-1]}$};
        \node [rectangle, minimum width=.5cm, minimum height=.75cm, text centered, text width=.75cm,  draw=black, fill=red!20] (hh2) at (-0.5, -6.5) {$\mathbf{S}_{[i]}$};
        \node [rectangle, minimum width=.5cm, minimum height=.75cm, text centered, text width=.75cm,  draw=black, fill=red!20] (hh3) at (1, -6.5) {$\mathbf{S}_{[i+1]}$};

        \node [rectangle, minimum width=.5cm, minimum height=.7cm, text centered, text width=.7cm,  draw=black, fill=orange!30] (hh4) at (3, -6.5) {$\mathbf{S}_{[i]}$};
        
        \node [rectangle, minimum width=.5cm, minimum height=1cm, text centered, text width=1cm,  draw=black, fill=blue!30] (h2) at (2, -1.5) {$\rmS_{[i]}$};
        \node [rectangle, minimum width=.5cm, minimum height=1cm, text centered, text width=1cm,  draw=black, fill=blue!30] (h3) at (6, -1.5) {$\rmS_{[i+1]}$};
        
        \node[draw, align=center,        draw=red!50,fill=red!20,rectangle split, 
        rectangle split horizontal,
        rectangle split parts=2, 
    %
    %
    ] (block4) at (4, 0){
        \nodepart{one} $\rmO^{\text{inter}}_{[i+1]}$ \nodepart{two} $\rmO^{\text{intra}}_{[i+1]}$ };
            
\node[draw,
        align=center,
        %
        draw=red!50,
        fill=red!20,
        rectangle split, 
        rectangle split horizontal,
        rectangle split parts=2, 
    %
    %
    ] (blockk4) at (5.5, -6.5){
        \nodepart{one} $\rmO^{\text{inter}}_{[i+1]}$ \nodepart{two} $\rmO^{\text{intra}}_{[i+1]}$
    };

    \node[draw,
        align=center,
        %
        draw=red!50,
        fill=red!20,
        rectangle split, 
        rectangle split horizontal,
        rectangle split parts=2, 
    %
    %
    ] (block2) at (0,0)
    {
        \nodepart{one} $\rmO^{\text{inter}}_{[i]}$ \nodepart{two} $\rmO^{\text{intra}}_{[i]}$
    };
    
    \node[BLOCK=3, fill=orange!30, draw=orange!80
              %
    %
    %
    ] (block) at (0,-3) {
        \nodepart{one} $\rmQ_{[i]}$ \nodepart{two} $\rmK_{[i]}$ \nodepart{three} $\rmV_{[i]}$
};

    \node[BLOCK=2, fill=orange!30, draw=orange!80
              %
    %
    %
    ] (blockk1) at (-2,-8) {
      \nodepart{one} $\rmK_{[i]}$ \nodepart{two} $\rmV_{[i]}$
};
    \node[BLOCK=2, fill=orange!30, draw=orange!80
              %
    %
    %
    ] (blockk2) at (0.5,-8) {
      \nodepart{one} $\rmK_{[i+1]}$ \nodepart{two} $\rmV_{[i+1]}$
};
    \node[BLOCK=3,  fill=orange!30, draw=orange!80
              %
    %
    %
    ] (block3) at (4,-3) {
        \nodepart{one} $\rmQ_{[i+1]}$ \nodepart{two} $\rmK_{[i+1]}$ \nodepart{three} $\rmV_{[i+1]}$
        
    };
    \node[BLOCK=3,  fill=orange!30, draw=orange!80
              %
    %
    %
    ] (blockk3) at (4.75,-8) {
        \nodepart{one} $\rmQ_{[i+1]}$ \nodepart{two} $\rmK_{[i+1]}$ \nodepart{three} $\rmV_{[i+1]}$
    };

    \node [rectangle, text centered, text width=.5cm, text height=.5cm, draw=black, fill=orange!30, label=right:{Load from HBM}] (node1) at (-2.5, -4.25){};
    \node [rectangle, text centered, text width=.5cm, text height=.5cm, draw=black, fill=red!20, label=right:Store to HBM] (node2) at (.9, -4.25){};
    \node [label=right:(a)] (labell) at (6, -3.3){};
        \node [label=right:(c)] (labelc) at (6, -8.6){};
            \node [label=right:(b)] (labelb) at (1, -8.6){};
    \node [rectangle, text centered, text width=.5cm, text height=.5cm, draw=black, fill=blue!30, label=right:On-chip construct] (node3) at (4, -4.25){};

            \draw [arrow] (process3) -- (block2.one south);
                \draw [arrow] (h2) -- (block4.one south);
            \draw [arrow] (block.one north) -- (block2.one south);
            \draw [arrow, dashed] (process3) -- (h2);
            \draw [arrow, dashed] (h2) -- (h3);
            \draw [arrow] (block.one north) -- (block2.two south);
                    \draw [arrow] (block3.one north) -- (block4.one south);
                                    \draw [arrow] (block3.two north) -- (block4.two south);
            \draw [arrow] (block3.three north) -- (block4.two south);
            \draw [arrow] (block3.three north) -- (h3);
            \draw [arrow] (block3.one north) -- (block4.two south);
            \draw [arrow] (block.two north) -- (block2.two south);
            \draw [arrow] (block.three north) -- (block2.two south);
            \draw [arrow] (block.two north) -- (h2);
            \draw [arrow] (block.three north) -- (h2);
            \draw [arrow, dashed] (hh1) -- (hh2);
            \draw [arrow, dashed] (hh2) -- (hh3);
            \draw [arrow] (blockk1.one north) -- (hh2);
            \draw [arrow] (blockk1.two north) -- (hh2);
            \draw [arrow] (blockk2.two north) -- (hh3);
            \draw [arrow] (blockk2.one north) -- (hh3);
            \draw [arrow] (hh4) -- (blockk4.one west);
            \draw [arrow] (blockk3.one north) -- (blockk4.one west);
                        \draw [arrow] (blockk3.one north) -- (blockk4.two south);
            \draw [arrow] (blockk3.two north) -- (blockk4.two south);
            \draw [arrow] (blockk3.three north) -- (blockk4.two south);
                        \draw [arrow, dashed] ($(hh1)-(1,0)$) -- (hh1);
                        \draw [arrow, dashed] (hh3) -- ($(hh3)+(1,0)$);
                        \draw [arrow, dashed] (h3) -- ($(h3)+(1.5,0)$);
                        \draw [arrow, dashed] ($(process3)-(1.5,0)$) -- (process3);


        \begin{scope}[on background layer]
    \node  [container1,fit= {(process3) (h2) (block) (block2) (h3) (block3) (labell)}, label=north:\large Sequential] (contain1) {};
                %
                %
                %
                %
                %
    \end{scope}
                        \begin{scope}[on background layer]
                \node  [container1,fit= { (hh1) (hh2) (hh3) (blockk1)  (blockk2) (labelb)}, label=north:\large Sequential] (contain2) {};
            \end{scope}
                        \begin{scope}[on background layer]
                \node  [container1,dcs, fit= { (hh4) (blockk3) (blockk4) (labelc)}, label={[label distance=0.25cm]north:\large Chunkwise Parallel}, fill=yellow!20] (contain3) {};
            \end{scope}

%
%
%
%
%
%

%
        \end{tikzpicture}   
}
        \vspace{-2mm}
        \caption{(a) \textsc{FlashLinearAttention} without materialization. This version is more memory-efficient. (b-c) \textsc{FlashLinearAttention} with materialization. This version enables sequence-level chunkwise parallelism.} \label{fig:chunk}
              
    \end{figure}